\section{Introduzione}
Questo capitolo esplicita i limiti informativi dell'\acrfull{adrr} rispetto al problema del \acrfull{dac} e dà conto, con trasparenza, anche dei tentativi non andati a buon fine. In particolare, l’elaborazione delle metriche temporali di Cap.~\ref{sec:metriche} calcolate sugli attraversamenti \gls{fir} (\S\ref{sec:flight_airspace}) non ha restituito segnali stabili o interpretabili in termini di configurazioni settoriali. La causa è di natura strutturale: \gls{adrr} rappresenta in modo accurato voli e tratte (\S\ref{sec:flights} e \S\ref{sec:flight_points}), ma non rende osservabili gli elementi centrali per il \gls{dac} delineati in Cap.~\ref{ch:introduzione-contesto} (\S\ref{subsec:settore-elementare}–\S\ref{subsec:settore-collassato}).

    Per superare tale limite, adottiamo un’ipotesi di lavoro allineata all’unica componente geometrica effettivamente osservabile in \gls{adrr}: le \gls{fir} sono assimilate a unità settoriali elementari. In concreto, una \gls{fir} è trattata come regione atomica (nessuna ulteriore sotto–settorizzazione osservabile). A partire da questa assunzione, costruiamo istanze coerenti con la geografia dei dati: le configurazioni ``elementari'' coincidono con singole \gls{fir}, mentre i ``collassati'' sono definiti come unioni connesse di più \gls{fir}, rispettando connettività e aderenza ai confini.

    %TODO citare i vincoli anche se per mancanza di dati non tutti i vincoli possono essere rispettati (settori collassati)
  Nel seguito si formalizza l’impiego delle \gls{fir} quali ``mattoni'' geometrici per la definizione di configurazioni elementari e collassate; si introducono i criteri di assegnazione della capacità e i vincoli geometrici; infine, si illustra l’integrazione di tali configurazioni nella formulazione del \gls{dac}.